\section{Rò rỉ tương tác}
\label{sec:ch11-ro-ri-tuong-tac}

\index{rò rỉ tương tác|(}%
Tính chất không mong muốn ấy có tên trong bài là
\tn{rò rỉ tương tác}{interaction leakage}, và bài phát biểu nó thẳng: các
phương pháp nối tầng, tự trong cấu trúc của chúng, không tách sạch được hiệu
ứng riêng của một đặc trưng khỏi hiệu ứng tương tác của nó với đặc trưng
khác~\autocite{p23metagame}. Muốn thấy được một phát biểu như vậy thì phải có
chỗ để so, nghĩa là phải có một trường hợp mà đáp án đúng đã biết trước. Bài
dựng đúng một trường hợp như thế.

Hàm được chọn là
%
\begin{equation}
  \label{eq:ch11-ham-vi-du}
  f(x) = x_1 + T, \qquad T := x_1 x_2^2,
\end{equation}
%
với baseline đặt tại gốc, $b = (0,0)$. Hàm này có đúng hai đặc trưng và tách
sẵn thành hai phần mà người đọc nhìn là biết: một phần chỉ phụ thuộc $x_1$, và
một phần chỉ khác không khi cả hai đặc trưng cùng khác không. Phần thứ hai là
tương tác, và sách đặt cho nó cái tên $T$ để khỏi phải viết lại. Bài báo viết
$I$ ở chỗ này, nhưng phụ lục~\ref{app:ky-hieu} đã gán $I$ cho
\tn{phép nhiễu}{perturbation} trong
định nghĩa infidelity của chương~\ref{ch:do-do-trung-thuc}, nên sách đổi tên
chứ không dùng lại ký hiệu.

Đáp án đúng ở đây không phải một quan sát về mô hình mà là một hệ quả của định
nghĩa. Bài định nghĩa \tn{hiệu ứng riêng}{pure individual effect} của đặc trưng
$i$ là attribution của $i$ khi chỉ mình nó có mặt, tức
$\varphi_{i \to i}(x,f) := \varphi_i(\{i\}; x, f)$: vẫn là phương pháp bậc nhất
chạy như thường, chỉ khác là chạy trên đầu vào đã bị che hết trừ đặc trưng $i$.
Với hàm ở phương trình~\ref{eq:ch11-ham-vi-du} và baseline tại gốc, che đặc
trưng thứ hai để lại hàm $x_1$, còn che đặc trưng thứ nhất để lại hàm hằng
bằng không. Nên mọi phương pháp bậc nhất trong bảng đều trả về $x_1$ cho đặc
trưng thứ nhất và $0$ cho đặc trưng thứ hai, bất kể chúng khác nhau thế nào ở
đầu vào đầy đủ; bài kiểm điều đó cho từng phương pháp trong phần chứng minh
phụ lục. Hai giá trị ấy bị ép bởi hàm và bởi baseline, không phải đo được từ
bên trong một mô hình nào. Bài phát biểu tiêu chuẩn dựa trên
chúng bằng đúng chữ mà cả Phần~IV đang bàn: một phương pháp trung thực phải
tách được hiệu ứng riêng của $x_1$ khỏi tương tác $x_1 x_2^2$, và cả serial
Shapley value lẫn integrated Hessians đều không làm được, vì chúng để các số
hạng tương tác rò sang các thành phần riêng~\autocite{p23metagame}.

Cách dựng ground truth ấy là cách chương~\ref{ch:chi-so-khong-do-do-trung-thuc}
đã đọc kỹ, chỉ khác đối tượng. Ở bài báo neo, đáp án đúng có được vì nhiệm vụ
được thiết kế sao cho chỉ có một đường đi tới câu trả lời; ở đây, đáp án đúng có
được vì hàm được viết ra sao cho hai phần của nó tách rời. Cái giá cũng cùng một
loại. Bảng~\ref{tab:ch11-tach-hieu-ung} nói về một đa thức hai biến với một quy
ước che đặc trưng đã cho sẵn, không nói về transformer mà mục~\ref{sec:ch11-ba-phep-do-va-ba-cai-nen}
sẽ tới; một phương pháp tách sạch ở đây chưa vì thế mà tách sạch ở kia.

\begin{table}[tbp]
  \centering
  \small
  \caption{Attribution bậc hai của ba nhóm phương pháp trên
  hàm ở phương trình~\ref{eq:ch11-ham-vi-du}, chép từ bảng phân tích của
  bài~\autocite{p23metagame}. Dòng đầu là của sách, tính từ định nghĩa hiệu ứng
  riêng của chính bài.}
  \label{tab:ch11-tach-hieu-ung}
  \begin{tabular}{@{}llcccc@{}}
    \toprule
    & & \multicolumn{2}{c}{Hiệu ứng riêng} & \multicolumn{2}{c}{Tương tác} \\
    \cmidrule(lr){3-4}\cmidrule(l){5-6}
    Nhóm & Phương pháp
      & $\varphi_{1\to1}$ & $\varphi_{2\to2}$
      & $\varphi_{2\to1}$ & $\varphi_{1\to2}$ \\
    \midrule
    \multicolumn{2}{@{}l}{\emph{Đáp án đúng}}
      & $x_1$ & $0$ & \multicolumn{2}{c}{\emph{tùy phương pháp bậc nhất}} \\
    \midrule
    Nối tầng & serial Shapley value
      & $x_1 + \tfrac{1}{4}T$ & $\tfrac{1}{4}T$
      & $\tfrac{1}{4}T$ & $\tfrac{1}{4}T$ \\
    & integrated Hessians
      & $x_1 + \tfrac{1}{9}T$ & $\tfrac{4}{9}T$
      & $\tfrac{2}{9}T$ & $\tfrac{2}{9}T$ \\
    \midrule
    Theo tập & biến đổi Möbius
      & $x_1$ & $0$ & \multicolumn{2}{c}{$T$, không có hướng} \\
    & STII
      & $x_1$ & $0$ & \multicolumn{2}{c}{$T$, không có hướng} \\
    & SOP
      & $x_1$ & $0$ & \multicolumn{2}{c}{$T$, không có hướng} \\
    \midrule
    Có hướng & Meta-G$\times$I
      & $x_1$ & $0$ & $T$ & $2T$ \\
    & Meta-IG
      & $x_1$ & $0$ & $\tfrac{1}{3}T$ & $\tfrac{2}{3}T$ \\
    & Meta-SV
      & $x_1$ & $0$ & $\tfrac{1}{2}T$ & $\tfrac{1}{2}T$ \\
    \bottomrule
  \end{tabular}
\end{table}

Hai ô trong bảng ấy đủ để thấy chỗ rò. Serial Shapley value trả về
$x_1 + \tfrac{1}{4}T$ cho hiệu ứng riêng của đặc trưng thứ nhất, tức một phần
tư tương tác đã bị cộng vào một con số lẽ ra chỉ được bằng $x_1$. Integrated
Hessians trả về $\tfrac{4}{9}T$ cho hiệu ứng riêng của đặc trưng thứ hai, mà
hiệu ứng riêng của đặc trưng ấy bằng không: bốn phần chín tương tác đã bị gán
cho một đặc trưng mà một mình nó không làm gì cả. Cả hai con số đều là đáp án
của một thủ tục chạy đúng như nó được định nghĩa, nên đây không phải lỗi cài
đặt mà là hệ quả của cấu trúc nối tầng.

Nhóm thứ hai trong bảng sửa được chỗ ấy nhưng đánh mất một thứ khác. Ba dòng
của nó là ba chỉ số chấm cả một tập đặc trưng cùng lúc: biến đổi Möbius, chỉ số
tương tác Shapley-Taylor mà bài viết tắt là STII, và SOP. Cả ba đều trả đúng
$x_1$ và $0$ cho hai hiệu ứng riêng, còn tương tác thì chúng để trong một ô duy nhất: chúng chấm cặp không có
thứ tự $\{1,2\}$, nên không có chỗ nào để phân biệt ảnh hưởng của đặc trưng thứ
hai lên attribution của đặc trưng thứ nhất với chiều ngược lại. Đó là chỗ mà
nhóm thứ ba, nhóm mà bài báo đề xuất, đi vào; và mục~\ref{sec:ch11-tro-choi-dung-tren-attribution}
dựng lại xây dựng ấy.
\index{rò rỉ tương tác|)}%
